\documentclass[a4paper,10pt]{amsart}

\pdfinfoomitdate=1
\pdftrailerid{}
\pdfsuppressptexinfo=15

\usepackage[T1]{fontenc}
\usepackage{lmodern}
\usepackage[margin=27mm]{geometry}
\usepackage{amsmath,amssymb,mathtools,booktabs,array,microtype,xcolor}
\usepackage[numbers,sort&compress]{natbib}
\usepackage[colorlinks=true,linkcolor=blue!55!black,citecolor=blue!55!black,urlcolor=blue!55!black]{hyperref}
\hypersetup{pdftitle={},pdfauthor={},pdfsubject={},pdfkeywords={}}

\newtheorem{theorem}{Theorem}
\newtheorem{lemma}[theorem]{Lemma}
\newtheorem{corollary}[theorem]{Corollary}
\theoremstyle{remark}
\newtheorem{remark}[theorem]{Remark}

\newcommand{\Sym}{\mathfrak S}
\newcommand{\std}{\operatorname{std}}
\newcommand{\CME}{\mathsf C}
\newcommand{\rtlmin}{\operatorname{rlmin}}
\newcommand{\stirfirst}[2]{\genfrac{[}{]}{0pt}{}{#1}{#2}}

\title[Cycle-maximum extraction]{Cycle-Maximum Extraction on Permutations:\\
Exact Images and Weighted Fibres}
\author{Anonymous}
\date{}

\begin{document}

\begin{abstract}
Write a permutation as disjoint cycles, order the cycle supports by their
least elements, read their greatest elements, and standardize the resulting
word.  Iterating this rule gives a rank-changing map on the disjoint union of
finite symmetric groups.  For a target $\sigma\in\Sym_m$, we prove that its
least possible source rank is
\[
 \mu(\sigma)=2m-\rtlmin(\sigma),
\]
and construct a source in every rank at least $\mu(\sigma)$.  The proof is an
explicit opener--closer schedule: precisely the right-to-left minima may use
one coordinate simultaneously as a cycle minimum and maximum.  Independently,
we resolve every target fibre as a sum over ordered set partitions, with a
block of size $b$ carrying weight $(b-1)!$.  The only recurrent permutations
are the identities, and every other step strictly lowers rank.  Static cycle
maxima, ordered-cycle conventions, and set-partition endpoint technology are
treated as prior inputs.  A deterministic audit executes $16{,}473{,}121$
exact assertions; it is counterexample pressure rather than proof.  No sharp
absorption clock is claimed.
\end{abstract}

\maketitle

\section{The map and its ownership boundary}\label{sec:map}

Let $\Sym_n$ be the permutations of $[n]$ in functional notation.  For
$\pi\in\Sym_n$, write $B_1,\ldots,B_m$ for the supports of its disjoint
cycles, ordered so that
\[
 \min B_1<\min B_2<\cdots<\min B_m.
\]
For a word of distinct integers, $\std$ replaces its smallest letter by $1$,
its next smallest by $2$, and so on.  Define the \emph{cycle-maximum
extraction map}
\begin{equation}\label{eq:map}
 \CME(\pi)=\std(\max B_1,\ldots,\max B_m)\in\Sym_m.
\end{equation}
On $\Sym_{\le N}=\bigsqcup_{1\le n\le N}\Sym_n$, this is a finite
self-map.  For example, the permutation with supports
\[
 \{1,3,5\},\quad\{2,4\},\quad\{6\}
\]
maps to $\std(5,4,6)=213$; cyclic order within a support is invisible to
\eqref{eq:map}.

The ingredients around \eqref{eq:map} are mature and receive no contribution
credit.  Chen--Deng--Du--Stanley--Yan fix block minima and maxima in their
crossing/nesting theory~\cite{ChenDengDuStanleyYan2005}; Rubey--Stump preserve
opener and closer configurations~\cite{RubeyStump2009}; Mongelli explicitly
orders cycles by increasing minima~\cite{Mongelli2012}; and
Andrews--Egge--Gawronski--Littlejohn use prescribed sets of cycle
maxima~\cite{AndrewsEggeGawronskiLittlejohn2011}.  We therefore claim neither
the endpoint language, the ordering convention, nor the elementary count of
cyclic orders.  The scope here is the exact inverse geometry of the literal
map \eqref{eq:map}: target-dependent image thresholds, all-rank sections, and
target-resolved weighted fibres.  The source search was bounded; a non-hit for
the exact map is not evidence of novelty, priority, or clearance.

For $\sigma=\sigma_1\cdots\sigma_m$, let $\rtlmin(\sigma)$ denote the number
of indices $i$ such that $\sigma_i<\sigma_j$ for all $j>i$.  Put
\begin{equation}\label{eq:mu}
 \mu(\sigma)=2m-\rtlmin(\sigma).
\end{equation}
For $n\ge m$, let $\mathcal P_n(\sigma)$ consist of the ordered set
partitions $(B_1,\ldots,B_m)$ of $[n]$ satisfying
\begin{equation}\label{eq:support-class}
 \min B_1<\cdots<\min B_m,
 \qquad \std(\max B_1,\ldots,\max B_m)=\sigma.
\end{equation}

\begin{theorem}[exact images, fibres, and recurrence]\label{thm:main}
For every $\sigma\in\Sym_m$ and $n\ge m$:
\begin{enumerate}
\item\label{it:image}
$\sigma\in\CME(\Sym_n)$ if and only if $n\ge\mu(\sigma)$.  The proof
constructs a deterministic right section in every admissible rank.
\item\label{it:fibre}
The complete target fibre is
\begin{equation}\label{eq:fibre}
 |\CME_n^{-1}(\sigma)|
 =\sum_{(B_1,\ldots,B_m)\in\mathcal P_n(\sigma)}
       \prod_{i=1}^m(|B_i|-1)!.
\end{equation}
In particular, the sum is zero precisely below the threshold in
\ref{it:image}.
\item\label{it:recurrent}
$\CME(\pi)=\pi$ holds exactly for identity permutations.  Every
nonidentity step strictly lowers rank, so the recurrent states of
$\Sym_{\le N}$ are $\mathrm{id}_1,\ldots,\mathrm{id}_N$.
\end{enumerate}
\end{theorem}

The image and fibre assertions are logically independent: the first resolves
existence and least rank through endpoints, whereas the second retains all
support choices and all cyclic orders.

\section{The endpoint schedule}\label{sec:image}

We first prove the lower bound in Theorem~\ref{thm:main}(\ref{it:image}).
Suppose $(B_1,\ldots,B_m)$ satisfies \eqref{eq:support-class}.  If $B_i$ is a
singleton, then for every $j>i$,
\[
 \max B_j\ge\min B_j>\min B_i=\max B_i.
\]
Thus $\sigma_i$ is a right-to-left minimum.  At most
$\rtlmin(\sigma)$ supports are singletons; all remaining supports use at least
two coordinates.  Hence
\[
 n=\sum_i|B_i|\ge\rtlmin(\sigma)+
 2\bigl(m-\rtlmin(\sigma)\bigr)=\mu(\sigma).
\]

The reverse implication is constructive.  Introduce formal opener and closer
chains
\begin{equation}\label{eq:chains}
 O_1<\cdots<O_m,\qquad K_1<\cdots<K_m,
\end{equation}
and pair $O_i$ with $K_{\sigma_i}$.  An opener records $\min B_i$, while its
paired closer records $\max B_i$.  A simultaneous event $S_i=O_i=K_{\sigma_i}$
will encode a singleton.

\begin{lemma}[greedy minimum schedule]\label{lem:schedule}
There is a precedence-respecting endpoint word of length
$2m-\rtlmin(\sigma)$ in which $O_i$ and $K_{\sigma_i}$ are simultaneous
exactly at the right-to-left minima of $\sigma$.
\end{lemma}

\begin{proof}
Let $i$ openers and $j$ closers have already been emitted, and continue while
$i<m$ or $j<m$.  If $i=m$, emit $K_{j+1}$; this is the forced final-closing
phase.  Otherwise $j<m$ as well, because an emitted closer requires its
owner to have opened.  Consider $O_{i+1}$ and $K_{j+1}$ and apply the
following deterministic rule:
\begin{enumerate}
\item if $i+1$ is a right-to-left-minimum position and
$\sigma_{i+1}=j+1$, emit the simultaneous event $S_{i+1}$;
\item otherwise, if the opener paired with $K_{j+1}$ has already appeared,
emit $K_{j+1}$;
\item otherwise emit $O_{i+1}$.
\end{enumerate}
The rule cannot become stuck: the boundary branch exhausts all remaining
closers after the last opener, and before then the third move is available
whenever the first two are not.

It remains to show that every designated simultaneous event is reached.  Let
$i+1$ be a right-to-left-minimum position and put $v=\sigma_{i+1}$.  Every
value smaller than $v$ occurs at a position at most $i$.  Its opener is
therefore available before $O_{i+1}$, and the closer-priority rule exhausts
$K_1,\ldots,K_{v-1}$ before opening $O_{i+1}$.  The next closer is then
$K_v$, so the first move applies.  No other position is declared
simultaneous.  Starting from $2m$ formal endpoints and making exactly
$\rtlmin(\sigma)$ identifications gives the stated length.
\end{proof}

Assign the integers $1,2,\ldots,\mu(\sigma)$ successively to the events in
Lemma~\ref{lem:schedule}.  For each $i$, take the coordinates at $O_i$ and
$K_{\sigma_i}$ as the endpoints of $B_i$; a simultaneous event gives a
singleton.  The opener chain orders the minima by $i$, and the closer chain
orders the maxima by their values, so \eqref{eq:support-class} holds.

This construction also gives every larger source rank.  Replacing a
simultaneous event by adjacent separate opener and closer events adds one
coordinate without changing either endpoint order.  After all simultaneous
events have been split, every support has two endpoints; further coordinates
may be inserted strictly between the opener and closer of one support.  Thus
there is a support family in $\mathcal P_n(\sigma)$ for every
$n\ge\mu(\sigma)$.  To obtain a permutation, put any one cycle on each
support---for instance, map its elements in increasing order cyclically.
This proves Theorem~\ref{thm:main}(\ref{it:image}), including the asserted
right sections.

For illustration, $\sigma=231$ has one right-to-left minimum, so
$\mu(231)=5$.  The supports of the permutation
\[
 \pi=(4,5,3,1,2)
\]
are $\{1,4\},\{2,5\},\{3\}$, and their maxima standardize to $231$.
The lower bound shows that rank five is best possible.

The scheduler is algorithmic as well as existential.  The inverse word
$\sigma^{-1}$ identifies which block owns the next closer; a right-to-left
scan marks the simultaneous positions.  The three-case rule in
Lemma~\ref{lem:schedule} then emits a minimum schedule in linear time after
these two scans.  Splits and interior insertions give a deterministic section
at any requested admissible rank.

\begin{corollary}[image census]\label{cor:image-count}
Let $\stirfirst{m}{r}$ be the unsigned Stirling number of the first kind.
Then
\begin{equation}\label{eq:image-count}
 |\CME(\Sym_n)|=
 \sum_{m=1}^n\ \sum_{r=\max(1,2m-n)}^m \stirfirst{m}{r}.
\end{equation}
At the minimum source rank of a target,
\begin{equation}\label{eq:min-fibre}
 |\CME_{\mu(\sigma)}^{-1}(\sigma)|
 =|\mathcal P_{\mu(\sigma)}(\sigma)|.
\end{equation}
\end{corollary}

\begin{proof}
Right-to-left minima on $\Sym_m$ have the unsigned first-kind Stirling
distribution: inserting the largest letter either creates a distinguished
record position or inserts into one of the existing gaps, giving the usual
Stirling recurrence.  The threshold $2m-r\le n$ now gives
\eqref{eq:image-count}.  This classical distribution receives no contribution
credit; the corollary merely specializes the target threshold.

If $n=\mu(\sigma)$, equality in the singleton lower bound forces precisely
the right-to-left-minimum blocks to have size one and every other block to
have size two.  Each factor in \eqref{eq:fibre} is then either $0!$ or $1!$,
so every feasible support family has weight one, proving
\eqref{eq:min-fibre}.
\end{proof}

For orientation, \eqref{eq:image-count} gives the first exact image sizes
shown below.
\begin{table}[ht]
\centering
\small
\caption{Exact literal profile used only as falsification pressure.}
\label{tab:profile}
\begin{tabular}{@{}lrrrrrrrrrr@{}}
\toprule
$n$ &1&2&3&4&5&6&7&8&9&10\\
\midrule
$|\CME(\Sym_n)|$&1&2&4&8&17&39&96&253&706&2074\\
\bottomrule
\end{tabular}
\end{table}

\section{Weighted fibres and dynamics}\label{sec:fibre}

\begin{proof}[Proof of Theorem~\ref{thm:main}(\ref{it:fibre})]
The support partition of a permutation is unique.  Fix
$(B_1,\ldots,B_m)\in\mathcal P_n(\sigma)$.  On a labelled set $B_i$ of size
$b_i$, there are $(b_i-1)!$ cyclic permutations: anchor one element and order
the remaining $b_i-1$ elements around it.  Cyclic orders on disjoint supports
are independent, giving $\prod_i(b_i-1)!$ sources with this support family.

Conversely, a support family satisfying \eqref{eq:support-class}, equipped
with one cyclic order on every block, determines a unique permutation, and
\eqref{eq:map} sends it to $\sigma$.  Distinct support families or cyclic
orders give distinct sources.  Summing the disjoint classes proves
\eqref{eq:fibre}.
\end{proof}

\begin{proof}[Proof of Theorem~\ref{thm:main}(\ref{it:recurrent})]
The rank of $\CME(\pi)$ is the number of cycles of $\pi$.  Equality with the
source rank $n$ occurs only when all $n$ cycles are singletons, that is, only
for $\pi=\mathrm{id}_n$.  Identities are fixed.  Every other orbit therefore
strictly decreases in positive integer rank until it reaches an identity, so
no other recurrent state exists.
\end{proof}

\section{Exact control and declarations}\label{sec:control}

The accompanying standard-library verifier independently constructs the
literal map through rank ten.  It checks $4{,}037{,}913$ source states, the
fixed/recurrent classification, $145{,}684$ target/rank image cells, $46{,}233$
minimum endpoint dynamic programs, $3{,}161$ explicit all-rank sections, and
$53{,}218$ fibre cells through source rank eight.  Restricted-growth words
generate the support side of \eqref{eq:fibre}; literal permutations generate
the other side.  The frozen run has $16{,}473{,}121$ exact assertions and ends
in \texttt{PASS}.  Enumeration is used only to look for counterexamples and
implementation errors; the proofs above are all-parameter arguments.

\paragraph{Limitations.}
The paper treats only the literal map \eqref{eq:map}.  It neither classifies
absorption times nor proves a sharp global clock.  Exact enumeration through
rank ten gives finite maximum tails
\(0,1,2,2,3,3,3,3,4,4\), suggesting a power-of-two pattern, but this is an
open computational observation only: no all-parameter bound, pointwise
clock, or global minimum-rank theorem for iterated preimages is asserted.
Static endpoint and cycle statistics are fully subtracted, and the bounded
direct-map search is not a novelty or priority claim.  No conclusion is
offered for other rules that choose representatives from cycles.

\paragraph{Data availability.}
All exact-control source code and its frozen text transcript are included in
the accompanying internal artifact.  They use no external data, runtime
network access, random sampling, or nonstandard Python package.

\paragraph{Ethics statement.}
This mathematical study involves no human participants, animals, personal
data, or field intervention.

\paragraph{Author contributions.}
The anonymous author is responsible for the definitions, proofs, software,
source audit, and manuscript.

\paragraph{Conflict of interest.}
The author declares no conflict of interest.

\paragraph{Funding.}
No external funding is declared.

\bibliographystyle{plainnat}
\bibliography{references}

\end{document}
